% =============================================================================
%  BATERIA DE PREGUNTAS PARA LA DEFENSA — FASE 4
%  Proyecto de Investigacion (ACIF3002) — Universidad Andres Bello
%  Daniel Alejandro Prado Correa
%
%  Documento de preparacion personal. No es un entregable de la rubrica.
%  Compilacion: pdfLaTeX
% =============================================================================

\documentclass[11pt,letterpaper,oneside]{article}

\usepackage[utf8]{inputenc}
\usepackage[T1]{fontenc}
\usepackage[spanish,es-nodecimaldot]{babel}

\usepackage{amssymb}
\let\Bbbk\relax
\usepackage{newtxtext}
\usepackage{newtxmath}

\usepackage[letterpaper,margin=2.2cm]{geometry}
\usepackage{amsmath}
\usepackage{enumitem}
\usepackage{xcolor}
\usepackage{mdframed}
\usepackage{titlesec}
\usepackage[hidelinks]{hyperref}

\definecolor{navy}{HTML}{173A5E}
\definecolor{rojo}{HTML}{C8102E}
\definecolor{gris}{HTML}{5A6672}
\definecolor{fondo}{HTML}{F4F7FB}

\titleformat{\section}{\large\bfseries\color{navy}}{\thesection.}{0.6em}{}
\titlespacing*{\section}{0pt}{1.4em}{0.7em}

\setlength{\parskip}{0.35em}
\setlength{\parindent}{0pt}

% --- entorno de pregunta -----------------------------------------------------
\newcounter{preg}
\newcommand{\pregunta}[1]{%
  \refstepcounter{preg}%
  \vspace{0.6em}%
  {\bfseries\color{rojo} P\thepreg.}\ {\bfseries #1}\par\vspace{0.15em}%
}
\newcommand{\resp}[1]{%
  \begingroup\color{gris}\small\hangindent=1.1em\hangafter=0%
  \textbf{$\rightarrow$}\ #1\par\endgroup\vspace{0.1em}%
}
\newcommand{\ojo}[1]{%
  \begingroup\color{navy}\small\itshape\hangindent=1.1em\hangafter=0%
  Nota: #1\par\endgroup%
}

\begin{document}

% =============================================================================
\begin{center}
{\LARGE\bfseries\color{navy} Batería de preguntas para la defensa}\\[0.4em]
{\large Proyecto de Investigación (ACIF3002) --- Fase 4}\\[0.3em]
{\color{gris} Daniel Alejandro Prado Correa \quad$\cdot$\quad Universidad Andrés Bello}
\end{center}

\vspace{0.8em}

\begin{mdframed}[backgroundcolor=fondo,linewidth=0pt,innertopmargin=8pt,
                 innerbottommargin=8pt,innerleftmargin=10pt,innerrightmargin=10pt]
\small
\textbf{Cómo usar este documento.} Las preguntas de la sección~1 se construyeron
directamente a partir de las observaciones que el docente guía formuló de manera
recurrente en las retroalimentaciones de las formativas y sumativas previas. Esas
observaciones son el mejor predictor disponible de lo que se preguntará en la
defensa: cada una de ellas ya fue corregida en los documentos finales, de modo
que la respuesta correcta consiste en \emph{señalar dónde} quedó resuelta. Las
secciones 2 a 5 cubren el terreno técnico y metodológico habitual de una comisión.

\smallskip
Las respuestas están redactadas en forma breve y hablada, no como texto de
informe. Practicar en voz alta: ninguna respuesta debería exceder los
\textbf{45 segundos}.
\end{mdframed}

% =============================================================================
\section{Preguntas derivadas de la retroalimentación del docente guía}

\pregunta{El informe está en normas APA y el artículo usa formato IEEE. ¿No es
una inconsistencia de criterio?}
\resp{No: son dos productos con normas distintas y ambas exigidas por la
rúbrica. El informe ABP sigue APA 7.\textsuperscript{a} edición mediante
\texttt{biblatex} con estilo \texttt{apa}; el artículo sigue el estilo numérico
IEEE mediante \texttt{IEEEtran}. Lo importante es que ambos productos comparten
un único archivo \texttt{referencias.bib}: el contenido bibliográfico es
literalmente el mismo archivo, solo cambia el estilo de salida. Por construcción
no pueden divergir.}

\pregunta{¿Dónde está el plan del proyecto y cómo verificó su cumplimiento?}
\resp{En la sección 7 del informe. El plan asocia a cada objetivo específico sus
actividades, su entregable auditable y su indicador de cumplimiento. La
verificación uno a uno está en la sección de conclusiones, y el Anexo~C rastrea
cada observación docente hasta la sección donde quedó incorporada.}

\pregunta{Usted reporta MSE y RMSE. ¿En qué unidades está cada uno?}
\resp{El MSE está en metros cúbicos a la sexta, porque es el cuadrado de un
volumen; el RMSE está en metros cúbicos, que es la unidad física interpretable.
Esa es exactamente la razón por la que la inferencia estadística la hago sobre el
RMSE y no sobre el MSE: el intervalo de confianza expresado en metros cúbicos
tiene sentido físico directo, el expresado en metros a la sexta no.}
\ojo{Esta observación apareció cuatro veces en las retroalimentaciones. Es la
pregunta más probable del bloque.}

\pregunta{¿Verificó el supuesto de normalidad antes de aplicar la prueba t?}
\resp{Sí, con Shapiro-Wilk, y el supuesto \textbf{se rechaza}: $p = 0{,}0233$.
Precisamente por eso el diseño metodológico incorporó desde el inicio el
contraste de Wilcoxon, que no exige normalidad y que confirma la conclusión con
$p = 8{,}88 \times 10^{-16}$. La conclusión se sostiene con independencia del
supuesto paramétrico. Además, con $M = 50$ el teorema central del límite da
robustez a la prueba t, y una $d$ de Cohen de 19 no se revierte por una
desviación distribucional moderada.}

\pregunta{¿Dónde está la matriz de revisión bibliográfica?}
\resp{En el Anexo~A del informe: doce fuentes caracterizadas por tipo, aporte al
estudio y DOI. Cada DOI fue verificado individualmente contra la fuente primaria.}

\pregunta{En la sección de ética, ¿qué riesgos identificó y cómo los mitigó?}
\resp{La sección está estructurada como riesgo seguido de mitigación, no como
declaración general. El riesgo principal es el uso de datos operacionales sujetos
a acuerdos de confidencialidad; la mitigación es trabajar con datos sintéticos
generados a partir de parámetros tomados de literatura técnica indexada, lo que
además permite conocer la verdad de referencia. Los demás riesgos --- integridad
de los resultados reportados, trazabilidad de las fuentes y reproducibilidad ---
están tratados con la misma estructura.}

\pregunta{En los avances anteriores el estado inicial aparecía como
$[15{,}0;\ 0{,}0]$ y en otro documento como $[10{,}0;\ 0{,}0]$. ¿Cuál es el
correcto?}
\resp{Detecté esa inconsistencia y la resolví de forma explícita. El estado
inicial real del sistema es $[15{,}0;\ -0{,}8]$ y el estado inicial del filtro es
$[15{,}0;\ 0{,}0]$: son cosas distintas y antes estaban confundidas. La
resolución está documentada en la sección 10 del informe.}

% =============================================================================
\section{Preguntas técnicas sobre el modelo y el estimador}

\pregunta{¿Por qué un modelo de dos estados y no uno solo?}
\resp{Porque con un solo estado no hay nada que reconstruir: si el estado es el
volumen y yo mido el volumen, el filtro solo suaviza. El segundo estado, la tasa
de llenado, es una variable que ningún instrumento mide, y es la que de verdad
importa operacionalmente: es la que anticipa el atoro. El modelo de dos estados
es lo que convierte el problema de suavizado en un problema de estimación.}

\pregunta{¿Cómo sabe que la tasa de llenado es reconstruible si no la mide?}
\resp{No lo asumo, lo demuestro. Construí la matriz de observabilidad apilando
$\mathbf{H}$ y $\mathbf{H}\mathbf{A}$, y su rango es 2, igual a la dimensión del
espacio de estados. Eso significa que la secuencia de mediciones de la única
variable observada determina unívocamente el estado completo. Es una condición
formal, verificada en código como parte del experimento.}

\pregunta{¿Qué es la forma de Joseph y por qué la usó?}
\resp{Es una manera algebraicamente equivalente pero numéricamente estable de
actualizar la matriz de covarianza del error. La forma simplificada puede perder
la simetría y la positividad definida por acumulación de error de redondeo, sobre
todo en corridas largas. La forma de Joseph preserva esas propiedades por
construcción. Con mil pasos por corrida y cincuenta corridas, esa decisión
importa.}

\pregunta{¿Por qué el Filtro de Kalman y no un filtro pasa-bajos o un promedio
móvil?}
\resp{Un promedio móvil no tiene modelo físico del proceso: introduce retardo de
fase, no pondera la confiabilidad relativa de la medición, y sobre todo no puede
entregar la tasa de llenado, porque solo opera sobre la señal que ya existe. El
Filtro de Kalman incorpora la ecuación de conservación de masa como restricción y
es óptimo en error cuadrático medio bajo los supuestos del modelo. La diferencia
no es de grado: el promedio móvil suaviza, el filtro reconstruye.}

\pregunta{¿Por qué recalibró el parámetro $q_{22}$ de la matriz de covarianza de
proceso?}
\resp{Porque el valor originalmente documentado, 0,05, es físicamente
implausible. Verifiqué numéricamente que sobre mil pasos produce un paseo
aleatorio de la velocidad del orden de $\pm 10$ metros cúbicos por segundo,
cuando la tasa nominal de extracción es 0,8: trece veces mayor que la magnitud
que pretende modelar. Lo recalibré a $1\times10^{-4}$ y documenté el ajuste de
forma explícita en el informe y en el README del código, en lugar de dejarlo
pasar.}

\pregunta{La ganancia de Kalman converge. ¿Qué implicancia práctica tiene eso?}
\resp{Es el comportamiento asintótico predicho por la ecuación de Riccati para
sistemas lineales invariantes en el tiempo. La implicancia práctica es
directa: una vez alcanzado el régimen permanente, en torno a los 60 segundos, la
ganancia podría precomputarse y quedar fija, lo que reduce el costo
computacional por iteración a un par de multiplicaciones matriciales. Para un
despliegue embebido eso es relevante.}

\pregunta{¿Qué ocurre si el sistema deja de ser lineal? Un atoro, un arranque.}
\resp{El modelo lineal deja de ser válido y el filtro degradaría. Es una
limitación que declaro explícitamente: el estudio se restringe al régimen de
operación continua. La vía de extensión es migrar a las variantes extendida o
\emph{unscented} del estimador, que es justamente el camino que siguen
Bárzaga-Martell y colaboradores en lixiviación de cobre. Está recogido como
segunda línea de continuidad.}

% =============================================================================
\section{Preguntas sobre validación estadística}

\pregunta{¿Por qué $M = 50$ corridas y no 500 o 1000?}
\resp{Porque 50 es suficiente para la pregunta que estoy respondiendo. El
intervalo de confianza al 95\,\% resultante tiene una amplitud de menos de
0,05 metros cúbicos sobre una diferencia media de 1,51: aumentar $M$ estrecharía
un intervalo que ya es estrecho, sin cambiar ninguna conclusión. El costo
computacional se justificaría si el efecto fuera marginal, y no lo es.}

\pregunta{¿Por qué una prueba pareada y no una de muestras independientes?}
\resp{Porque las observaciones están pareadas por construcción: en cada corrida,
la lectura cruda del SCADA y la estimación del filtro se evalúan sobre la misma
realización del ruido y la misma trayectoria verdadera. Tratarlas como
independientes ignoraría esa correlación y reduciría la potencia de la prueba
sin ninguna ganancia.}

\pregunta{Una $d$ de Cohen de 19 parece irreal. ¿No es un artefacto?}
\resp{Es alta porque la varianza \emph{entre corridas} es pequeña comparada con
la diferencia de medias: el efecto del filtro es muy consistente de corrida en
corrida. No es un artefacto, es lo que uno esperaría de una comparación
controlada donde la única fuente de variación es la semilla del ruido. Por eso
reporto también el intervalo de confianza y el contraste no paramétrico, que son
menos sensibles a esa escala.}

\pregunta{Los datos son sintéticos. ¿Eso no invalida los resultados?}
\resp{Delimita el alcance, no invalida el resultado. De hecho es una condición
necesaria del diseño: para calcular el error cuadrático medio hay que conocer la
verdad de referencia, y en una faena real esa verdad no existe, porque es
precisamente lo que no se puede medir. Lo que sí reconozco como limitación es que
no hay validación bajo discrepancia genuina de modelo, y esa es la primera línea
de continuidad del trabajo.}

\pregunta{¿Cómo garantiza que sus resultados son reproducibles?}
\resp{Con semilla pseudoaleatoria fija y flujos independientes generados con
\texttt{SeedSequence.spawn}, lo que da corridas estadísticamente independientes y
a la vez exactamente repetibles. Y lo verifiqué: dos ejecuciones sucesivas e
independientes produjeron archivos de salida idénticos byte a byte, comprobado
con suma MD5 sobre los cuatro registros CSV. Es un hecho verificado, no una
declaración de intención.}

% =============================================================================
\section{Preguntas sobre alcance, aporte y proyección}

\pregunta{Si el Filtro de Kalman existe desde 1960, ¿cuál es el aporte de este
trabajo?}
\resp{Ninguno algorítmico, y lo digo explícitamente en el informe. El aporte está
en tres planos. Arquitectónico: formular la estimación estocástica como
componente de diseño de un gemelo digital minero, no como instrumento de control
ni como preprocesamiento, que es la capa que las revisiones de Qu y de Ebad
identifican como desatendida. Empírico: entregar evidencia cuantificada de que la
diferencia entre una sombra digital y un gemelo digital es medible. Y
metodológico: un artefacto reproducible byte a byte con verificación de supuestos
y análisis de robustez, en un área donde la validación efectiva está señalada
como carencia.}

\pregunta{¿Esto es investigación o es desarrollo de software?}
\resp{Es investigación aplicada con un artefacto como instrumento. La pregunta de
investigación no es «se puede programar un filtro de Kalman» --- eso está
resuelto hace décadas --- sino si una capa de estimación estocástica, formulada
como componente arquitectónico, reconstruye de forma confiable el estado interno
de un activo parcialmente observable. Eso se responde con evidencia estadística,
y el software es el medio para producirla.}

\pregunta{¿En qué se diferencia su módulo de un preprocesamiento cualquiera del
\emph{pipeline} de datos?}
\resp{En que un preprocesamiento opera solo sobre la señal que existe: la
suaviza, la interpola, le quita valores atípicos. Mi módulo produce una variable
que no existe en el sistema. Esa es la distinción que quiero dejar planteada: no
se trata solo de limpiar los datos, sino de reconstruirlos. Y eso justifica
tratarlo como componente arquitectónico y no como un accesorio del
\emph{pipeline}.}

\pregunta{¿Cómo se integraría esto a un gemelo digital real?}
\resp{El módulo está deliberadamente desacoplado: recibe la serie de mediciones y
las matrices del sistema, y entrega el estado estimado con su covarianza. No
depende de base de datos, protocolo industrial ni interfaz. Se ubicaría entre la
capa de adquisición y cualquier consumidor analítico. Como evidencia de ese
desacoplamiento, apliqué el mismo núcleo sin modificar una línea a otros tres
subsistemas del activo: potencia del motor, temperatura del estator y desgaste
del revestimiento.}

\pregunta{¿Cuál es el siguiente paso?}
\resp{Validación con telemetría real de faena, que es condición necesaria para
cualquier despliegue y permitiría medir el desempeño bajo discrepancia genuina de
modelo. Después, extensión a régimen no lineal y sintonización adaptativa de las
covarianzas a partir de la secuencia de innovaciones, que es lo que sugiere
directamente el hallazgo sobre la razón $\mathbf{Q}/\mathbf{R}$.}

% =============================================================================
\section{Preguntas incómodas}

\pregunta{Los resultados de sus avances anteriores no coinciden con estos. ¿Por
qué cambiaron?}
\resp{Porque el código original no implementaba el modelo que los documentos
describían. Implementaba un modelo escalar de un solo estado, no el modelo de dos
estados que estaba documentado, y además no fijaba semilla, por lo que sus
resultados no eran reproducibles. Reconstruí el \emph{pipeline} conforme al
modelo documentado y reporto lo que efectivamente arroja. La diferencia en las
cifras es consecuencia de la corrección, no del método, y prefiero declararlo
antes que sostener un número que no puedo reproducir.}
\ojo{Responder esto de frente y sin rodeos es una fortaleza, no una debilidad.
Una comisión valora más a quien corrige su propio trabajo que a quien defiende un
número frágil.}

\pregunta{¿Verificó que todas sus referencias existen?}
\resp{Sí, una por una contra la fuente primaria. Retiré las que no pude
verificar y las reemplacé por fuentes con DOI comprobado. También corregí una
referencia que estaba atribuida a la revista equivocada. Las doce referencias
finales están en el Anexo~A con su DOI, y la declaración de originalidad del
informe deja constancia de este procedimiento.}

\pregunta{¿Utilizó herramientas de inteligencia artificial en este trabajo?}
\resp{Responder con transparencia según lo que efectivamente se usó y conforme a
la política institucional. El argumento de fondo que sí corresponde sostener es:
los resultados no dependen de ninguna declaración, provienen de código
ejecutable con semilla fija que cualquiera puede correr y verificar, y puedo
explicar y defender cada decisión de modelado, cada parámetro y cada prueba
estadística de este trabajo.}

\pregunta{Su reducción del MSE es del 89\,\%, pero la del RMSE es solo del
68\,\%. ¿No está inflando la cifra?}
\resp{No, son la misma cosa expresada en escalas distintas: el RMSE es la raíz
del MSE, así que una reducción del 89\,\% en la escala cuadrática equivale
exactamente a una del 68\,\% en la escala lineal. Reporto ambas justamente para
que no haya ambigüedad, y hago la inferencia estadística sobre el RMSE, que es la
métrica en unidades físicas.}
\ojo{Verificación: $1 - \sqrt{1 - 0{,}8957} = 0{,}6771$. Coherente con el
67,80\,\% reportado, dentro del redondeo de las medias por corrida.}

% =============================================================================
\section{Frases de rescate}

Si una pregunta no tiene respuesta a mano, estas fórmulas permiten mantener el
control sin improvisar contenido:

\begin{itemize}[leftmargin=1.3em,itemsep=0.35em]
\item «Esa pregunta excede el alcance que declaré para el estudio. La recojo como
línea de continuidad, y le comento cómo la abordaría.»
\item «No tengo esa cifra exacta en la cabeza, pero está en la sección X del
informe y puedo mostrársela.»
\item «Déjeme distinguir dos cosas que creo que se están mezclando en la
pregunta.»
\item «Eso es una limitación real del trabajo y está declarada como tal. Lo que
sí puedo sostener con la evidencia que tengo es \dots»
\end{itemize}

\vspace{1.2em}
\begin{mdframed}[backgroundcolor=fondo,linewidth=0pt,innertopmargin=8pt,
                 innerbottommargin=8pt,innerleftmargin=10pt,innerrightmargin=10pt]
\small
\textbf{Cifras que conviene tener memorizadas.}\quad
Reducción del MSE del volumen: 89,57\,\% \quad$\cdot$\quad
Reducción del error en la tasa latente: 98,17\,\% \quad$\cdot$\quad
$t(49) = 135{,}24$ \quad$\cdot$\quad
$d$ de Cohen $= 19{,}13$ \quad$\cdot$\quad
Shapiro-Wilk $p = 0{,}0233$ (se rechaza) \quad$\cdot$\quad
Wilcoxon $p = 8{,}88\times10^{-16}$ \quad$\cdot$\quad
Rango de la matriz de observabilidad $= 2$ \quad$\cdot$\quad
$M = 50$ corridas, $N = 1000$ pasos, semilla 42.
\end{mdframed}

\end{document}
